\begin{itemize}
    \color{red}
    \item what is ACTS and where does it come from
    \item what is its purpose and goals
    \item how is the project structured
\end{itemize}

ACTS stands for A Common Tracking Software and is a community-driven project which aims to provide experiment independent reconstruction algorithms written in modern C++. Apart from that ACTS can be used for algorithm prototyping and early performance studies. The project is open source and can be found on GitHub.

The project is structured in a modular way, where each component is responsible for a specific task, such as magnetic field, material, geometry, propagation, track finding, track fitting, and vertexing. These components interact with each other in various ways.

ACTS consists of two main parts: the Core library and the Examples framework. The Core library is a toolbox and meant to be used by experiments to build their own reconstruction framework using tools provided by ACTS. The Examples framework demonstrates how to use the Core library and provides a set of tools to run algorithms on a simulated detector. The Examples framework can be used for educational purposes, testing and early performance studies but it is generally disencouraged to use it for production workflows.

The reason for this separation is that a reconstruction framework is highly experiment specific and should be tailored to the specific needs of the experiment. The Core library is meant to be used as a building block for such a framework. The Examples framework is meant to be used as a starting point for new users to get familiar with the Core library.

\section{Components and their Interactions}

ACTS can be thought of a collection of components which interact with each other to build higher and higher level of abstractions. The most basic components are the magnetic field, the geometry and the material. The magnetic field can be queried for the field strength at a given point in space. There are a few implementations available with the simplest one being a constant field. The geometry component deals with surfaces, layers and volumes which then build up the tracking geometry. The material component abstracts material properties and provides utilities to estimate interactions with the material.

The propagator component is responsible for moving particles along their trajectory. It can be used to extrapolate track parameters from one surface to another. The propagator uses the magnetic field to calculate the trajectory of the particle. The geometry is used to check for intersections with the tracking geometry and the material is used to estimate energy loss and multiple scattering.

Ultimately we build algorithms on top of these components. Some of them depend on the propagation, like classical track finding and fitting algorithms, while others just depend on the geometry, like track seeding algorithms.

Around these components there are a few more definitions and concepts which do not map directly necessarily to single components in the code. Two of these are units and constants and the track parameterization.

\section{Units and Constants}

To make sense of physical quantities, ACTS uses a pre-defined set of units. They can be found in \texttt{Acts/Definitions/Units.hpp}.

Most importantly, lengths are measured in millimeters and energies in GeV. The speed of light is set to $1$ which means that the time unit is also in millimeters. Seconds can be obtained by dividing the time with pre-defined constants which are related to the SI units.

\section{Track Parameterization}

Track parameters describe the state of a particle at a given point in space and/or time. This includes the position, time, direction and charge/momentum of the particle. There are two kinds of track parameters in ACTS: bound parameters and free parameters. Bound track parameters are expressed with a local position at a given surface and their direction with two spherical coordinate angles, while free track parameters are expressed with a global cartesian position and a normalized direction vector.

\begin{equation*}
    \vec x_{local} = \left(l_0, l_1, \phi, \theta, q/p, t\right)^T
\end{equation*}

\begin{equation*}
    \vec x_{global} = \left(g_0, g_1, g_2, t, d_0, d_1, d_2, q/p\right)^T
\end{equation*}

\begin{itemize}
    \color{red}
    \item how does one convert from bound to free
    \item mention covariance
    \item angles and their value range
\end{itemize}

In ACTS \texttt{BoundTrackParameters} contain the bound parameter vector, the surface and optionally a covariance matrix. \texttt{FreeTrackParameters} contain the free parameter vector and optionally a covariance matrix.

Track parameters are the input for the propagation and can be extrapolated along the trajectory. This is ultimately the basis for the track fitting.

Note that the ACTS bound track parameters are sort of standard for collider experiments. This comes with two caveats: the bound parametrization does not make a lot of sense for foward exteriments and there are poles in the parametrization which can lead to numerical instabilities. This can be avoided by rotating the tracking geometry away from the poles.

The track parametrization is ultimately a choice. There is not single best parametrization. Since ACTS aims to be experiment independent a more general solution would be preferable but was not reached yet. The parametrization is deeply baked into ACTS which makes it hard to change. At the same time an abstraction for the track parameters might bring a compute performance penalty.

Track parameters should be unique and "well-behaved" in the sense that they should not have singularities or discontinuities. This is not always the case with the ACTS track parameters. Additionally, for numerical stability, they should live on a similar scale which also depends on the units and geometry.

\section{Plugins}

\begin{itemize}
    \color{red}
    \item what are plugins
    \item which plugins are available
\end{itemize}

\section{The ACTS Examples Framework}

\begin{itemize}
    \color{red}
    \item what is the difference to the core library
    \item how is it structured
    \item components: geometry, magnetic field, sequencer, whiteboard, algorithms, reader, writer
\end{itemize}

As mentioned before, the Examples framework excercised the use of the Core library to build a reconstruction framework. It can be used for educational purposes, testing and early performance studies. It is generally discouraged to use it for production workflows. As such, the Examples are not subject to the same level of scrutiny as the Core library. They are meant to be a starting point for new users to get familiar with the Core library.

The goal of the Examples is to test common workflows and to enable monitoring of our algorithms on a higher level. Testing within the Core library cannot include a detailed simulation and reconstruction chains. The Examples framework provides a way to test the Core library in a more realistic environment.

The Examples consist of various detectors, such as a generic telescope detector and the Open Data Detector, a sequencer, which executes reconstruction algorihtms in a given order, a whiteboard, which is a data store for the algorithms, a reader and writer, which read and write event data and summaries to and from disk, and a set of algorithms, such as track finding, track fitting and vertexing.

\section{The Open Data Detector}

\begin{itemize}
    \color{red}
    \item dd4hep
    \item evolution of the track ML detector
    \item what is its purpose
    \item full silicon detector
    \item digitization and resolution
    \item support structures
    \item layout and geometry, subsystems, eta coverage
    \item material scans
    \item expected hits over eta/phi
    \item expected resolution
\end{itemize}

\section{Tracking on Accelerators with traccc and detray}

\begin{itemize}
    \color{red}
    \item very brief - link to other peoples work
\end{itemize}
